\inicializarSubseccion{Diferencias entre for y while}
\tab Es verdad que tanto el \texttt{for} como el \texttt{while} comparten una estructura de grupo isomorfa, sin embargo, en contextos de programación, se suelen usar en diferentes situaciones. Las diferencias entre el ciclo for y while es que el \texttt{for} se usa cuando sabes el número de repeticiones desde el inicio, mientras que en el \texttt{while} no conoces el número de iteraciones. For es mejor para estructuras finitas como arreglos o listas mientras que el while es más útil para situaciones donde no es fija la condición de parada